% TOP_PROBLEM: {"id": 30006122, "title": "Arthur-Langlands Reciprocity", "category_id": 1, "category": {"id": 1, "name": "number_theory", "display_name": "Number Theory", "description": "Properties of integers, prime numbers, Diophantine equations.", "slug": "number-theory", "order_index": 1, "created_at": "2026-07-31T15:26:25.670Z"}, "status": "open"}
% FIRST_POSTED: 2026-10-06
% ATTEMPT_STATUS: unresolved
% OPENAI_CLAIM: {"schema_version":1,"source_repo":"https://github.com/openai/math","source_commit":"adc7f1241b42e322a6451854ab7e4b4c146bf78a","release_date":"2026-10-06","match":"partial","resolution":"partial","independently_reviewed":false,"second_pass":"n/a","families":[{"family":"014","title":"Restricted geometric Langlands, global Arthur enhancements, and generic Ramanujan","subject":"Number theory","manuscripts":[{"title":"Constructible tame Hecke eigensheaves in positive characteristic","pdf_path":"preprints/Constructible-tame-Hecke-eigensheaves-in-positive-characteristic-October-5-2026/constructible-tame-hecke-eigensheaves-positive-characteristic.pdf","readme_path":"preprints/Constructible-tame-Hecke-eigensheaves-in-positive-characteristic-October-5-2026/README.md","date":"2026-10-05","theorem_ref":"Theorem (label thm:main), preprints/Constructible-tame-Hecke-eigensheaves-in-positive-characteristic-October-5-2026/build/manuscript/introduction.tex"},{"title":"Global Arthur Enhancements of Cuspidal Excursion Parameters","pdf_path":"preprints/Global-Arthur-Enhancements-of-Cuspidal-Excursion-Parameters-October-5-2026/manuscript.pdf","readme_path":"preprints/Global-Arthur-Enhancements-of-Cuspidal-Excursion-Parameters-October-5-2026/README.md","date":"2026-10-05","theorem_ref":"Theorem (label thm:main), preprints/Global-Arthur-Enhancements-of-Cuspidal-Excursion-Parameters-October-5-2026/build/manuscript.tex"},{"title":"Ramanujan-Arthur Decompositions of Cuspidal Functions at Full Finite Level","pdf_path":"preprints/Ramanujan-Arthur-Decompositions-of-Cuspidal-Functions-at-Full-Finite-Level-September-24-2026/paper.pdf","readme_path":"preprints/Ramanujan-Arthur-Decompositions-of-Cuspidal-Functions-at-Full-Finite-Level-September-24-2026/README.md","date":"2026-09-24","theorem_ref":"Theorem (label main:decomposition), preprints/Ramanujan-Arthur-Decompositions-of-Cuspidal-Functions-at-Full-Finite-Level-September-24-2026/build/sections/01-introduction.tex"},{"title":"Rationality of the Canonical Unramified Arthur Filtration","pdf_path":"preprints/Rationality-of-the-Canonical-Unramified-Arthur-Filtration-September-24-2026/paper.pdf","readme_path":"preprints/Rationality-of-the-Canonical-Unramified-Arthur-Filtration-September-24-2026/README.md","date":"2026-09-24","theorem_ref":"Theorem (label intro:main), preprints/Rationality-of-the-Canonical-Unramified-Arthur-Filtration-September-24-2026/build/sections/01-introduction.tex"},{"title":"The Restricted Geometric Langlands Equivalence in Positive Characteristic","pdf_path":"preprints/The-Restricted-Geometric-Langlands-Equivalence-in-Positive-Characteristic-September-24-2026/paper.pdf","readme_path":"preprints/The-Restricted-Geometric-Langlands-Equivalence-in-Positive-Characteristic-September-24-2026/README.md","date":"2026-09-24","theorem_ref":"Theorem (label thm:full-support), preprints/The-Restricted-Geometric-Langlands-Equivalence-in-Positive-Characteristic-September-24-2026/build/sections/01-introduction.tex"}],"lean":null,"reasoning_trace":null}]}
\documentclass[11pt]{article}
\usepackage[margin=25mm]{geometry}
\usepackage{fontspec}
\setmainfont{Latin Modern Roman}
\usepackage{amsmath,amssymb,mathtools}
\usepackage{unicode-math}
\setmathfont{Latin Modern Math}
\usepackage{enumitem}
\usepackage{xurl,hyperref,bookmark}
\hypersetup{colorlinks=true,urlcolor=blue}
\setcounter{secnumdepth}{0}
\setlength{\parindent}{0pt}
\setlength{\parskip}{5pt}
\setlength{\emergencystretch}{3em}
\begin{document}
\section*{Arthur-Langlands Reciprocity}

\subsection{OpenAI's claim}
OpenAI presents the following claims relevant to this question, with manuscript numbers following paper order in Sources. In manuscript 1 (family 014; 2026-10-05), the asserted result is that tame regular-unipotent, geometrically dense parameters admit nonzero constructible perverse Hecke eigensheaves with one Borel-level point for the specified simply connected groups and characteristic hypotheses; In manuscript 2 (family 014; 2026-10-05), the asserted result is that assuming the stated finite-level decomposition, each occurring cuspidal excursion parameter for a split semisimple function-field group has a global Arthur enhancement with one algebraic $\mathrm{SL}_2$ and a commuting Weil action, including inertia; In manuscript 3 (family 014; 2026-09-24), the asserted result is that cuspidal functions for split semisimple groups over function fields decompose by dual nilpotent orbits at every full finite level; for adjoint absolutely simple groups, one generic unramified component forces temperedness at every unramified place; In manuscript 4 (family 014; 2026-09-24), the asserted result is that the canonical unramified Arthur filtration on all finitely supported automorphic functions is rational under the restricted theory's characteristic hypotheses for split connected semisimple groups; In manuscript 5 (family 014; 2026-09-24), the asserted result is that restricted geometric Langlands has full spectral support in positive characteristic in the two stated characteristic regimes, with additional very-good and Weyl-order conditions over arbitrary algebraically closed fields.

\subsection{Outline of the argument}
For manuscript 1 (family 014; 2026-10-05), the parameter is glued through a separating node to an unramified family. Two specializations and transverse Hecke cuts detect a nonzero flag-level object; unipotent monodromy on the enhancement torus permits descent to ordinary Borel level. For manuscript 2 (family 014; 2026-10-05), shtuka cohomology supplies a finite-type cyclic support carrying full excursion data. Frobenius-compatible coherent models of Iwahori categories at several places identify an open-orbit condition, which assembles a single commuting Weil and $\mathrm{SL}_2$ pair.

For manuscript 3 (family 014; 2026-09-24), iwahori translation eigencharacters are analyzed before taking the Weyl quotient. Frobenius traces of central kernels and Springer tests bound the extreme Bernstein exponent, while trace identities synchronize the governing nilpotent orbit across places and coefficient embeddings. For manuscript 4 (family 014; 2026-09-24), a bilinear detector tests supported spectral local cohomology against finite evaluations. Orbit distributions, compact Weil lifts, and noncancellation separate the nilpotent supports, and rational descent then includes the noncuspidal functions.

For manuscript 5 (family 014; 2026-09-24), a hypothetical missing component is transported along a trait to a geometric specialization. Levi reduction and transverse Hecke modifications provide cotangent tests whose nonvanishing prevents that spectral summand from being annihilated.

\subsection{Scope relative to this problem}
The linked constructions give specified function-field and positive-characteristic branches: full-support equivalence, rational filtrations, selected cuspidal decompositions, generic temperedness, and tame eigensheaves. The conditional global enhancement also keeps its decomposition premise. They do not supply the general discrete-spectrum, packet, multiplicity, and reciprocity description for arbitrary reductive groups and fields asked here.

\subsection{Formal verification}
No Lean formalization is listed at this snapshot for manuscript 1 (family 014; 2026-10-05), manuscript 2 (family 014; 2026-10-05), manuscript 3 (family 014; 2026-09-24), manuscript 4 (family 014; 2026-09-24), manuscript 5 (family 014; 2026-09-24).

This site has not executed the Lean code. Described statement coverage is not an independent proof check.

\subsection{Status caveat}
The manuscripts were produced by an internal OpenAI model. OpenAI cautions that results without formalization can contain mistakes. This site has not independently checked the proofs; the label records the scope of OpenAI's claim.

\subsection{Sources}
\begin{itemize}
\item OpenAI, \emph{Constructible tame Hecke eigensheaves in positive characteristic}, 2026-10-05; citation key \texttt{OAI:Constructible-tame-Hecke-eigensheaves-in-positive-characteristic-October-5-2026}. \href{https://github.com/openai/math/blob/adc7f1241b42e322a6451854ab7e4b4c146bf78a/preprints/Constructible-tame-Hecke-eigensheaves-in-positive-characteristic-October-5-2026/constructible-tame-hecke-eigensheaves-positive-characteristic.pdf}{PDF (GitHub)}; \href{https://github.com/openai/math/raw/adc7f1241b42e322a6451854ab7e4b4c146bf78a/preprints/Constructible-tame-Hecke-eigensheaves-in-positive-characteristic-October-5-2026/constructible-tame-hecke-eigensheaves-positive-characteristic.pdf}{PDF (direct)}; \href{https://github.com/openai/math/blob/adc7f1241b42e322a6451854ab7e4b4c146bf78a/preprints/Constructible-tame-Hecke-eigensheaves-in-positive-characteristic-October-5-2026/README.md}{Manuscript page and citation}; \href{https://github.com/openai/math/blob/main/preprints/Constructible-tame-Hecke-eigensheaves-in-positive-characteristic-October-5-2026/constructible-tame-hecke-eigensheaves-positive-characteristic.pdf}{Latest version}.
\item OpenAI, \emph{Global Arthur Enhancements of Cuspidal Excursion Parameters}, 2026-10-05; citation key \texttt{OAI:Global-Arthur-Enhancements-of-Cuspidal-Excursion-Parameters-October-5-2026}. \href{https://github.com/openai/math/blob/adc7f1241b42e322a6451854ab7e4b4c146bf78a/preprints/Global-Arthur-Enhancements-of-Cuspidal-Excursion-Parameters-October-5-2026/manuscript.pdf}{PDF (GitHub)}; \href{https://github.com/openai/math/raw/adc7f1241b42e322a6451854ab7e4b4c146bf78a/preprints/Global-Arthur-Enhancements-of-Cuspidal-Excursion-Parameters-October-5-2026/manuscript.pdf}{PDF (direct)}; \href{https://github.com/openai/math/blob/adc7f1241b42e322a6451854ab7e4b4c146bf78a/preprints/Global-Arthur-Enhancements-of-Cuspidal-Excursion-Parameters-October-5-2026/README.md}{Manuscript page and citation}; \href{https://github.com/openai/math/blob/main/preprints/Global-Arthur-Enhancements-of-Cuspidal-Excursion-Parameters-October-5-2026/manuscript.pdf}{Latest version}.
\item OpenAI, \emph{Ramanujan-Arthur Decompositions of Cuspidal Functions at Full Finite Level}, 2026-09-24; citation key \texttt{OAI:Ramanujan-Arthur-Decompositions-of-Cuspidal-Functions-at-Full-Finite-Level-September-24-2026}. \href{https://github.com/openai/math/blob/adc7f1241b42e322a6451854ab7e4b4c146bf78a/preprints/Ramanujan-Arthur-Decompositions-of-Cuspidal-Functions-at-Full-Finite-Level-September-24-2026/paper.pdf}{PDF (GitHub)}; \href{https://github.com/openai/math/raw/adc7f1241b42e322a6451854ab7e4b4c146bf78a/preprints/Ramanujan-Arthur-Decompositions-of-Cuspidal-Functions-at-Full-Finite-Level-September-24-2026/paper.pdf}{PDF (direct)}; \href{https://github.com/openai/math/blob/adc7f1241b42e322a6451854ab7e4b4c146bf78a/preprints/Ramanujan-Arthur-Decompositions-of-Cuspidal-Functions-at-Full-Finite-Level-September-24-2026/README.md}{Manuscript page and citation}; \href{https://github.com/openai/math/blob/main/preprints/Ramanujan-Arthur-Decompositions-of-Cuspidal-Functions-at-Full-Finite-Level-September-24-2026/paper.pdf}{Latest version}.
\item OpenAI, \emph{Rationality of the Canonical Unramified Arthur Filtration}, 2026-09-24; citation key \texttt{OAI:Rationality-of-the-Canonical-Unramified-Arthur-Filtration-September-24-2026}. \href{https://github.com/openai/math/blob/adc7f1241b42e322a6451854ab7e4b4c146bf78a/preprints/Rationality-of-the-Canonical-Unramified-Arthur-Filtration-September-24-2026/paper.pdf}{PDF (GitHub)}; \href{https://github.com/openai/math/raw/adc7f1241b42e322a6451854ab7e4b4c146bf78a/preprints/Rationality-of-the-Canonical-Unramified-Arthur-Filtration-September-24-2026/paper.pdf}{PDF (direct)}; \href{https://github.com/openai/math/blob/adc7f1241b42e322a6451854ab7e4b4c146bf78a/preprints/Rationality-of-the-Canonical-Unramified-Arthur-Filtration-September-24-2026/README.md}{Manuscript page and citation}; \href{https://github.com/openai/math/blob/main/preprints/Rationality-of-the-Canonical-Unramified-Arthur-Filtration-September-24-2026/paper.pdf}{Latest version}.
\item OpenAI, \emph{The Restricted Geometric Langlands Equivalence in Positive Characteristic}, 2026-09-24; citation key \texttt{OAI:The-Restricted-Geometric-Langlands-Equivalence-in-Positive-Characteristic-September-24-2026}. \href{https://github.com/openai/math/blob/adc7f1241b42e322a6451854ab7e4b4c146bf78a/preprints/The-Restricted-Geometric-Langlands-Equivalence-in-Positive-Characteristic-September-24-2026/paper.pdf}{PDF (GitHub)}; \href{https://github.com/openai/math/raw/adc7f1241b42e322a6451854ab7e4b4c146bf78a/preprints/The-Restricted-Geometric-Langlands-Equivalence-in-Positive-Characteristic-September-24-2026/paper.pdf}{PDF (direct)}; \href{https://github.com/openai/math/blob/adc7f1241b42e322a6451854ab7e4b4c146bf78a/preprints/The-Restricted-Geometric-Langlands-Equivalence-in-Positive-Characteristic-September-24-2026/README.md}{Manuscript page and citation}; \href{https://github.com/openai/math/blob/main/preprints/The-Restricted-Geometric-Langlands-Equivalence-in-Positive-Characteristic-September-24-2026/paper.pdf}{Latest version}.
\item \href{https://github.com/openai/math/blob/adc7f1241b42e322a6451854ab7e4b4c146bf78a/CONTENTS.md}{OpenAI catalogue entry}.
\end{itemize}
\end{document}
